\section{What the repository proves, and what this audit establishes}
\label{sec:audit}

\subsection{The endpoint is Dirichlet GRH}

We inspected the pinned source, not just the repository title.
The file \file{grh/src/Statements.lean} imports only Mathlib and states
the propositions proved by \file{grh/src/Main.lean}. The five exports are
listed below. In the general Dirichlet statement, \(N\ne0\), all characters
modulo \(N\) are allowed, and the pole of the trivial character at \(s=1\)
is excluded explicitly. The zero-on-line statement concerns
\(0<\Re s<1\). The separate zeta statement excludes its usual trivial
zeros. These are substantive statements, including imprimitive
characters, with no custom redefinition of GRH hiding in the
endpoint file~\cite{repo}.

\begin{center}
\small
\begin{tabularx}{\textwidth}{@{}p{76mm}Y@{}}
\toprule
Export in \file{Main.lean} & Mathematical scope\\
\midrule
\file{riemannHypothesis} & Mathlib's proposition \file{RiemannHypothesis}.\\
\file{riemannZeta_eq_zero_re} & Every zeta zero apart from the stated exclusions
has real part \(1/2\).\\
\file{LFunction_ne_zero_of_half_lt_re} &
Dirichlet \(L(\chi,s)\ne0\) for \(\Re s>1/2\), apart from the pole case.\\
\file{LFunction_eq_zero_re} &
Dirichlet zeros in the open critical strip lie on the critical line.\\
\file{LFunction_eq_zero_re_of_isPrimitive} &
The corresponding statement restricted to primitive characters.\\
\bottomrule
\end{tabularx}
\end{center}

This immediately supplies RH for abelian Dedekind zeta functions by their
Dirichlet factorization. It is not an export of arbitrary nonabelian
Dedekind, Hecke-over-number-fields, Artin, or automorphic GRH.
That scope is exactly why an \(A_5\) splitting field is a meaningful
further target.

\subsection{Reproducibility and trust boundary}

The readable development uses Lean 4.35.0 and Mathlib revision
\file{56a005855ea17e0b0958f73ec801ae140b0a1679}.
At the pinned commit, our source inventory finds 803 readable-edition
Lean files under \file{grh/src}, totaling 477,124 physical lines including
comments, and 794 compact-edition files totaling 112,827 lines.
The accompanying script records hashes of key files and the declaration
locations used in this report.

The repository and independent checking pages report successful builds,
kernel replay or independent checking, endpoint comparison, and only
\file{propext}, \file{Classical.choice}, and \file{Quot.sound} for the
exported theorems~\cite{repo,sutherland-check,sonoda-check,rubinstein-check}.
This report accepts the mathematical endpoint specified by the task and
uses those records as external verification evidence. It does not claim
to have repeated them. No Lean toolchain was installed in this workspace,
and no full \texttt{lake build}, \texttt{\#print axioms}, or independent
kernel replay was run here.

Our limited lexical scan strips comments and strings and finds no
standalone \file{sorry}, \file{admit}, \file{axiom}, or \file{unsafe}
tokens in either edition's \file{src} tree. That is useful source hygiene,
not a kernel audit or a proof about dependency trust. An expected
placeholder in a challenge file outside those source trees is not evidence
of an unfinished theorem in the compiled endpoint.

The three mathematical developments below were separately cross-checked
for their substantive algebraic and analytic steps. The finite numerical
certificates were executed. The derived extensions are written proofs
in this report; none is represented as a new Lean-checked theorem.

\section{Which ideas can actually be transplanted?}
\label{sec:mechanism}

\subsection{The arithmetic comparison that produces the saving}

The proof's distinctive content is the construction and comparison of
arithmetic lattices of forms. The final Landau argument is an analytic
consumer of an arithmetic estimate. It cannot on its own generate the
estimate for another \(L\)-function.

At a schematic level, short integral forms with a controlled index give
an upper bound for a normalized covolume. A sufficiently large independent
family makes the useful prime remainder visible with a substantial
coefficient. Long-form estimates from theta chains, congruences, and
Selmer-controlled twists give a lower bound with the same leading term.
In the source's notation the comparison has the shape
\[
 M-\text{controlled error}
 \ \le\ \frac{V+2\delta_{\mathrm{lat}}}{d}
 \ \le\ M-\frac nd\,mR+\text{controlled error}.
\]
Here \(M\) is the common main term, \(d\) the target dimension, \(n\) the
spanning dimension, \(V\) a covolume term, and
\(\delta_{\mathrm{lat}}\) an index contribution. The subtraction of the
common \(M\) leaves an upper bound on \(mR\).
This formula is explanatory notation, not a replacement for the exact
Lean inequalities~\cite{repo,alpoge-remark}.

\begin{center}
\small
\begin{tabularx}{\textwidth}{@{}>{\raggedright\arraybackslash}p{28mm}YY@{}}
\toprule
Stage & Source location & Role in a transfer\\
\midrule
Short and enough &
\file{New/Reduction/ShortAndEnough.lean} &
Integral short vectors, index control, and enough independent vectors.\\
\addlinespace
Local-to-global lattice &
\file{New/Reduction/Lattices/TwistedLattice.lean} &
Product formula converts locally chosen twists and denominator bounds
into a global covolume inequality.\\
\addlinespace
Long forms &
\file{New/Determinant/Twist/LongForms.lean} &
Arithmetic lower bound whose main term matches the short-form bound.\\
\addlinespace
Cancellation and scale &
\file{New/Reduction/MainInequality.lean};
\file{New/Reduction/OneSidedBound.lean} &
\file{star_inequality} gives
\(m\,\mathrm{totalRemainder}\le fa(\log b)^5\), then \(m\asymp\sqrt{x}\).\\
\addlinespace
Zeros &
\file{New/Reduction/Zeros/Criterion.lean};
\file{Lib/Analytic/Landau.lean} &
A one-sided bound and appropriate real residue signs exclude poles of
a transform, hence unwanted zeros.\\
\bottomrule
\end{tabularx}
\end{center}
All paths in this table are relative to \file{grh/src}.

In particular, an intermediate bound in
\file{eventually_primeRemainder_natCast_le} is
\[
 R_q(x)\le \bigl(64\varphi(q)+3\bigr)x^{-1/2+6\varepsilon}
 \qquad\text{eventually}.
\]
The choice of \(\varepsilon\), the fixed modulus, the auxiliary scales,
and the source's prime-number-theorem input lead to the final requested
one-sided exponent. The square-root scale comes from the length of
the arithmetic construction; it is not supplied by an abstract
zero-free principle.

The actual prime coefficients involve congruence conditions and a
quadratic factor of the form
\[
 \Lambda(n)\bigl(1-\chi_K(n)\bigr)
             \mathbf 1_{\,n\equiv1\pmod q}.
\]
The source also selects auxiliary quadratic twists to prevent cancellation
among residues in a prescribed box. A replacement target must justify
both its coefficients and its noncancellation input. Relabeling
\(L(\chi,s)\) as a general automorphic \(L\)-function does neither.

\subsection{The local quantifiers matter}

The twisted-lattice theorem permits the local twist to depend on the
rational prime and on the embedding into a \(p\)-adic algebraic closure.
Its product-formula conclusion does not assert that a single global
character works at every finite place. Similarly, selecting factors
before taking their products or gcds can change a divisibility statement.
The earlier reports already emphasize this quantifier issue; it must
survive every attempted transfer.

The torsion theorem in Section~\ref{sec:torsion-transfer} has a different,
particularly useful uniformity: at one fixed prime it chooses a finite
grid before seeing the family of polynomials or analytic functions.
That uniformity makes an infinite-family extension possible. A static
import-closure scan of its source file reaches 103 local modules,
58,528 lines, and 3,388,790 bytes, with Mathlib as the only external
import name. This is the current syntactic import closure, not a
minimal mathematical dependency set. It contains no final RH endpoint
and makes this theorem a plausible standalone extraction target.

\subsection{What the five preceding reports already cover}
\label{sec:prior}

All five ZIPs in the selected \file{docs/incoming} snapshot were inspected.
Their hashes and exact internal report paths are recorded in
\file{evidence/prior_reports_manifest.json}. The table summarizes their
claimed completed arguments and their explicit gaps; enumerations
recorded in those packets were not all rerun in this analysis~\cite{prior}.

\begin{center}
\footnotesize
\begin{tabularx}{\textwidth}{@{}p{49mm}YY@{}}
\toprule
Archive & Already developed there & Boundary relevant here\\
\midrule
\file{grh_transfer_analysis.zip} &
Determinant sandwich, scalar Landau reduction, conductor-23 cubic target,
and seven small algebraic Lean lemmas. &
No principal-class prime bound and no nonabelian RH conclusion.\\
\addlinespace
\file{grh_transfer_analysis_2026-11-25.zip} &
Higher-order pole criterion, sparse witness families, Tate-inertia
obstruction, and product/gcd quantifier analysis. &
No arithmetic construction of the sparse family; no Vandiver or
\(K\)-group vanishing theorem.\\
\addlinespace
\file{grh-transfer.zip} &
A reported 229-field abelian sufficient-condition classification,
small-prime representations, split-prime criteria, and upper limsup. &
Reported enumeration not rerun here; an equivalent prime inequality
is not a proved inequality.\\
\addlinespace
\file{grh-lean-transfers.zip} &
A reported 26-field quintic CM classification, genus containment,
weighted races, and \(S_3\) detector analysis. &
No \(S_3\) zero theorem; algebraic classifications do not extend the
theta construction.\\
\addlinespace
\file{grh_transfers.zip} &
Partial exponent-six census, mod-4 window positivity, eta-product
variance, and a conductor-23 sign equivalence. &
The full exponent-six census remains unproved; it explicitly asks for
intermediate genus-field improvements.\\
\bottomrule
\end{tabularx}
\end{center}

The present report adds the analytic and infinite-family torsion theorem,
the \(A_5\) construction with its exact coefficients and variance, and the
intermediate genus-field improvement. It does not relabel the previous
scalar Landau theorem, the previously observed Tate-inertia obstruction,
or the existence of simultaneous nonvanishing results as new discoveries.

For comparison, combining RH with the established nonnegativity of the
de Bruijn--Newman constant gives \(\Lambda=0\) immediately. That is an
important consequence but not a transfer of the internal method, and it
is excluded from this report's new findings~\cite{rodgers-tao}.
The same restraint applies to familiar effective-prime, class-number,
and conditional algorithm consequences unless an additional argument
substantially improves the preceding analysis.

